\chapter{An Example Chapter}
\label{ch:example}

Use ordinary paragraphs for exposition. Start by explaining why a new object
is useful, then move to its formal definition. A new conceptual step deserves
a new paragraph, while formal statements belong in dedicated environments.

\section{Definitions and statements}

\begin{definition}[Subset]
  \label{def:subset}
  A set $A$ is a \term{subset} of a set $B$ if every element of $A$ also
  belongs to $B$. In this case, write $A \subseteq B$.
\end{definition}

Follow a definition with a brief interpretation. Here, $A \subseteq B$ means
that the set $A$ is entirely contained in $B$.

\begin{notation}
  The symbols $\N$, $\Z$, $\Q$, $\R$, and $\C$ denote the natural, integer,
  rational, real, and complex numbers, respectively.
\end{notation}

\begin{proposition}[Transitivity of inclusion]
  \label{prop:subset-transitive}
  If $A \subseteq B$ and $B \subseteq C$, then $A \subseteq C$.
\end{proposition}

\begin{proof}
  Choose any $x \in A$. Since $A \subseteq B$, we have $x \in B$. Since
  $B \subseteq C$, it follows that $x \in C$. Therefore every element of
  $A$ belongs to $C$.
\end{proof}

Use labels for reliable cross-references; see \cref{def:subset} and
\cref{prop:subset-transitive}.

\section{Problems and solutions}

\begin{problem}[Intersection and union]
  \label{prob:intersection-union}
  Prove that for any sets $A$ and $B$,
  \[
    A \cap B \subseteq A \cup B.
  \]
\end{problem}

\begin{solution}
  Let $x \in A \cap B$. Then $x \in A$ and $x \in B$. In particular, at
  least one of the conditions $x \in A$ or $x \in B$ holds, so
  $x \in A \cup B$.
\end{solution}

\begin{example}
  If $A=\{1,2\}$ and $B=\{2,3\}$, then
  $A\cap B=\{2\}$ and $A\cup B=\{1,2,3\}$.
\end{example}

\begin{remark}
  The argument uses only the definitions of intersection and union, not the
  specific elements of the sets.
\end{remark}

\begin{summary}
  Definitions answer ``What is it?'' Statements answer ``What is always
  true?'' Proofs explain why a statement is true. Problems stay separate from
  their solutions so that each condition is easy to find during review.
\end{summary}
